\documentclass[../main.tex]{subfiles}
\graphicspath{{\subfix{../images/}}}
\begin{document}
\subsection*{Answers - Turning problems into quadratics (page \pageref{quadratics})}
\label{Quadratics answers}
\begin{spacing}{1.5}
\begin{enumerate}
    \item
    \((2^2)^x+2^x-24=0\)\\
    \((2^x)^2+2^x-24=0\)\\
    Making the substitution \(u=2^x\)\\
    \(u^2+2u-24=0\)\\
    \(u=4.42, -5.42\)\\
    \(2^x\) can never be negative so we can ignore the -5.42 solution.\\
    \(2^x=4.42\)\\
    \(\ln{2^x}=\ln{4.42}\)\\
    \(x\ln{2}=\ln{4.42}\)\\
    \(x=\frac{\ln{4.42}}{\ln{2}}\)\\
    \(x=2.14\) (1dp)
    \item 
    Rearrange to \(9^x-6^x-4^x=0\)\\
    We need a constant so divide through by the lowest term.\\
    \(
    \frac{9^x}{4^x}-\frac{6^x}{4^x}-\frac{4^x}{4^x}=0\\
    (\frac{9}{4})^x-(\frac{6}{4})^x-1=0\\
    ((\frac{3}{2})^2)^x-(\frac{3}{2})^x-1=0\\
    ((\frac{3}{2})^x)^2-(\frac{3}{2})^x-1=0\\ 
    \)
    Use the substitution \(u=(\frac{3}{2})^x\)\\
    \(
    u^2-u-1=0\\
    u-1.618, -0.618\\
    \)
    \((\frac{3}{2})^x\) can never be negative so we ignore -0.618.\\
    \(
    (\frac{3}{2})^x=1.618\\
    \ln{(\frac{3}{2})^x}=\ln{1.618}\\
    x\ln(\frac{3}{2})=\ln{1.6.18}\\
    x=\frac{\ln{1.618}}{\ln{\frac{3}{2}}}\\
    x=1.187
    \)
    \item 
    We need a constant so divide through by the lowest term.\\
    \(
    8(\frac{9^x}{4^x})+3(\frac{6^x}{4^x})-81=0\\
    8(\frac{9}{4})^x+3(\frac{6}{4})^x-81=0\\
    8((\frac{3}{2})^2)^x+3(\frac{3}{2})^x-81=0\\
    8((\frac{3}{2})^x)^2+3(\frac{3}{2})^x-81=0\\
    \)
    Use the substitution \(u=(\frac{3}{2})^x\)\\
    \(8u^2+3u-81=0\)\\
    \(u=3,-3.375\)\\
    Since \((\frac{3}{2})^x\) can never be negative, we can ignore the -3.375 solution.\\
    \((\frac{3}{2})^x=3\)\\
    \(\ln{(\frac{3}{2})^x}=\ln{3}\)\\
    \(x\ln{(\frac{3}{2})}=\ln{3}\)\\
    \(x=\frac{\ln{3}}{\ln{\frac{3}{2}}}\)=2.71
    \item 
    We need a constant so divide through by the lowest term.\\
    \(
    (\frac{25^x}{9^x})+2(\frac{15^x}{9^x})-24=0\\
    (\frac{25}{9})^x+2(\frac{15}{9})^x-24=0\\
    ((\frac{5}{3})^2)^x+2(\frac{5}{3})^x-24=0\\
    ((\frac{5}{3})^x)^2+2(\frac{5}{3})^x-24=0\\
    \)
    Use the substitution \(u=(\frac{5}{3})^x\)\\
    \(u^2+2u-24=0\)\\
    \(u=4,-6\)\\
    Since \((\frac{5}{3})^x\) can never be negative, we can ignore the -6 solution.\\
    \((\frac{5}{3})^x=4\)\\
    \(\ln{(\frac{5}{3})^x}=\ln{4}\)\\
    \(x\ln{(\frac{5}{3})}=\ln{4}\)\\
    \(x=\frac{\ln{4}}{\ln{\frac{5}{3}}}=2.714\)

    \item
    Multiply through by \(e^{x}\):\\
    \(
    e^{2x} - 20 = e^{x}\\
    {(e^{x})}^2 - 20 = 20e^{x}\\
    {(e^{x})}^2 + e^{x} - 20 = 0\\
    e^{x} = -5, 4\\\)
    Can't be negative, therefore:\\
    \(
    e^{x}=4\\
    x = \ln{4}
    \)

    \item
    \(
    3^{2x+1} - 82\times 3^{x} +27 = 0\\
    3^{1}\times 3^{2x} - 82\times 3^{x} +27 = 0\\
    3^{1}\times {(3^{x})}^2 - 82\times 3^{x} +27 = 0\\
    u = 3^{x}\\
    3u^{2} - 82u + 27 = 0\\
    u = \frac{1}{3}, 27\\
    3^{x} = \frac{1}{3} \rightarrow x = -1\\
    3^{x} = 27 \rightarrow x = 3
    x = -1, 3
    \)

    \item 
    \(
    x^{\frac{1}{4}} +2x^{-\frac{1}{4}}=3\\
    {(x^{\frac{1}{4}})}^2 + 2 = 3x^{\frac{1}{4}}\\
    {(x^{\frac{1}{4}})}^2 - 3x^{\frac{1}{4}} + 2 = 0\\
    u = x^{\frac{1}{4}}\\
    u^2 - 3u + 2 = 0\\
    u = 1, 2\\
    x^{\frac{1}{4}} = 1 \rightarrow x = 1\\
    x^{\frac{1}{4}} = 2 \rightarrow x = 2^4 = 16\\
    x = 1, 16
    \)

    \item 
    \(x^2 - 5x + 6 + \frac{5}{x} + \frac{1}{x^2} = 0\)
    Group the terms with same powers:
    \(
    x^2 + \frac{1}{x^2} - 5x + \frac{5}{x} + 6 = 0\\
    x^2 + \frac{1}{x^2} - 5(x + \frac{1}{x}) + 6 = 0
    \)
    Use the substitution \(u = x + \frac{1}{x}\)
    \(
    {(x + \frac{1}{x})}^2 - 2 - 5(x + \frac{1}{x}) + 6 = 0\\
    {(x + \frac{1}{x})}^2 - 5(x + \frac{1}{x}) + 4 = 0\\
    u^2 - 5u + 4 = 0\\
    u = 1, 4
    \)
    Reverse substitution and solve:\\
    \(
    x + \frac{1}{x} = 1\\
    x^2 + 1 = x\\
    x^2 - x + 1 = 0\\
    x = \frac{1 \pm i\sqrt{3}}{2}\\
    x + \frac{1}{x} = 4\\
    x^2 + 1 = 4x\\
    x^2 - 4x + 1 = 0\\
    x = \frac{4 \pm 2\sqrt{3}}{2} = 2 \pm \sqrt{3}
    \)
    Therefore, \(x = \frac{1 + i\sqrt{3}}{2}, \frac{1 - i\sqrt{3}}{2}, 2 + \sqrt{3}, 2 - \sqrt{3}\)
\end{enumerate}


\end{spacing}
\end{document}